\documentclass[11pt,a4paper]{article}
\usepackage[margin=23mm]{geometry}
\usepackage[T1]{fontenc}
\usepackage{lmodern}
\usepackage{amsmath,amssymb}
\usepackage{algorithm}
\usepackage{algpseudocode}
\usepackage{float}
\usepackage{hyperref}
\hypersetup{hidelinks}
\renewcommand{\thealgorithm}{4.\arabic{algorithm}}
\algrenewcommand\algorithmicrequire{\textbf{Input:}}
\algrenewcommand\algorithmicensure{\textbf{Output:}}
\setlength{\parindent}{0pt}
\setlength{\parskip}{6pt}
\title{Chapter 4 Implementation Algorithms}
\author{}
\date{}

\begin{document}
\maketitle
These algorithms abstract the retained implementation. They do not introduce
duplicate-aware repartitioning or new model training. Algorithm 4.12 describes
the legacy speech pilot, not the current complete chatbot evaluation.
\begin{algorithm}[H]
\caption{Load prepared text data and retain the experiment partitions}
\label{alg:data}
\begin{algorithmic}[1]
\Require Dataset mode $d$, source files, seed $s=42$
\Ensure Training frame $T$, test frame $E$, class mapping $L$
\If{$d=\text{Symptom2Disease}$}
  \State Read the CSV; exclude unnamed index columns
  \State Keep text and label; drop missing or empty values; strip whitespace
  \State Assign a source row identifier to each usable record
  \State $(T,E)\gets\Call{RandomSplit}{\text{records},0.80,0.20,s}$
  \State Do not stratify the outer split
  \State Check the expected 960/240 sizes and disjoint source row IDs
\Else
  \State Read the supplied cleaned training and test CSVs separately
  \For{$F\in\{T,E\}$}
    \State Check diagnosis and cleaned-text columns exist
    \State Drop missing entries; convert to strings; strip whitespace
    \State Remove empty text or labels; reset row indices
  \EndFor
  \State Retain the supplied partition assignments
\EndIf
\State $L\gets\Call{EncodeLabelVocabulary}{\text{labels in }T\text{ and }E}$
\State Map each target label to its integer class identifier
\State Report row counts and class count
\State \Return $(T,E,L)$
\end{algorithmic}
\end{algorithm}
No text deduplication or group-aware splitting occurs in this procedure.
The preliminary BioBERT internal validation step is described separately in the chapter.

\newpage
\begin{algorithm}[H]
\caption{Train and evaluate TF-IDF with Logistic Regression}
\label{alg:tfidf}
\begin{algorithmic}[1]
\Require Frames $T,E$, class mapping $L$, TF-IDF and solver settings $\theta$
\Ensure Fitted pipeline $M$, predictions $\hat y$, training duration $t$
\State Configure word unigrams and bigrams
\State Enable lowercasing, English stop words and sublinear term frequency
\State Set feature cap and minimum document frequency from $\theta$
\State Configure Logistic Regression solver, iteration limit and seed
\State $M\gets\Call{Pipeline}{\text{TF-IDF},\text{Logistic Regression}}$
\State Start the training timer
\State Fit $M$ on the text and encoded labels in $T$
\State $t\gets$ elapsed training duration
\State $\hat y\gets\Call{Predict}{M,E.\text{text}}$
\State Calculate accuracy, macro and weighted metrics
\State Export the class report using the documented label convention
\State Save the fitted vectoriser and classifier together with joblib
\State \Return $(M,\hat y,t)$
\end{algorithmic}
\end{algorithm}
The vectoriser is fitted on training records only. The preliminary feature cap
is 20,000; the expanded cap is 50,000. The expanded implementation uses SAGA.

\newpage
\begin{algorithm}[H]
\caption{Fine-tune an expanded-dataset transformer classifier}
\label{alg:transformer}
\begin{algorithmic}[1]
\Require Checkpoint $q$, prepared frames $T,E$, labels $L$, settings $\theta$
\Ensure Saved classifier, tokeniser, predictions and trainer history
\State Load the tokeniser associated with $q$
\State Construct label-to-ID and ID-to-label mappings from $L$
\State Load the sequence-classification model with $|L|$ output classes
\State Initialise the task-specific classification head for this label space
\For{$F\in\{T,E\}$}
  \State Tokenise $F.\text{text}$ using truncation at 128 tokens
  \State Retain encoded targets and remove the raw text field
\EndFor
\State Configure learning rate, batch sizes, epochs, weight decay and seed
\State Configure padding through the tokeniser's data collator
\State Configure mixed precision only as supported by the runtime settings
\State Construct Trainer with $T$ as the training dataset
\State Start timer and run the configured training operation
\State Record runtime and any checkpoint-resumption provenance
\State $Z\gets\Call{TrainerPredict}{E}$
\State $\hat y_i\gets\arg\max_k Z_{ik}$ for each test record $i$
\State Save model weights, tokeniser, label mapping and trainer history
\State Export predictions and documented evaluation metrics
\State \Return saved artifacts and $\hat y$
\end{algorithmic}
\end{algorithm}
This algorithm represents the common expanded training structure. The preliminary
BioBERT notebook selects an epoch on 864/96 records and then trains a fresh model
on all 960 outer-training records. Saved runs may use different checkpoint cadences.

\newpage
\begin{algorithm}[H]
\caption{Enhance a decoded mono 16 kHz waveform}
\label{alg:audio}
\begin{algorithmic}[1]
\Require Decoded waveform $x$, sample rate $f_s=16{,}000$
\Ensure Enhanced waveform $x'$, preprocessing metadata $A$
\State Replace non-finite values; reject empty or effectively silent input
\State Record original sample count and duration
\State $\mu\gets\Call{Mean}{x}$
\State $x_c\gets x-\mu$
\State $(x_t,b)\gets\Call{TrimBoundaries}{x_c,30\text{ dB}}$
\If{$x_t$ is empty}
  \State $x_t\gets x_c$; retain the original centred interval
\EndIf
\State $x_d\gets\Call{StationaryNoiseReduce}{x_t,f_s,0.5}$
\State Replace any non-finite values in $x_d$
\State $a\gets\max|x_d|$
\If{$a>0$}
  \State $x_d\gets0.95x_d/a$
\EndIf
\State $x'\gets\Call{ValidateFloat32Audio}{x_d}$
\State $A\gets$ durations, removed offset, trim interval and enhancement settings
\State \Return $(x',A)$
\end{algorithmic}
\end{algorithm}
Decoding and resampling occur before this function. The original-audio comparison
still uses mono 16 kHz audio; it skips enhancement, not format standardisation.

\newpage
\begin{algorithm}[H]
\caption{Transcribe audio with the reusable Whisper runtime}
\label{alg:whisper}
\begin{algorithmic}[1]
\Require Waveform $x$, language English, chunk length 25 s, selected device
\Ensure Automatic unedited transcript $h$
\State Load and cache Whisper processor and model if not already loaded
\State Set model to evaluation mode
\State Split $x$ into consecutive chunks of at most $25f_s$ samples
\State $H\gets$ empty ordered list
\For{each chunk $c$ in recording order}
  \State Compute acoustic features and an attention mask with the processor
  \State Move available input tensors to the selected device
  \State Generate transcription tokens in inference-only mode
  \State Decode tokens, excluding special tokens; strip outer whitespace
  \If{decoded chunk text is non-empty}
    \State Append the chunk text to $H$
  \EndIf
\EndFor
\State $h\gets\Call{JoinInOrder}{H}$
\If{$h$ is empty}
  \State Raise an empty-transcription error
\EndIf
\State \Return $h$ for transcript confirmation or automatic evaluation
\end{algorithmic}
\end{algorithm}
The chunks do not overlap. Human transcript correction is a separate assistance
condition and must not replace the raw hypothesis in automatic WER/CER scoring.

\newpage
\begin{algorithm}[H]
\caption{Extract canonical symptoms and their assertion states}
\label{alg:extract}
\begin{algorithmic}[1]
\Require Narrative $n$, eligible canonical phrases and aliases $V$
\Ensure Grouped evidence $E$, current-symptom text $b$
\State $u\gets\Call{MatchingNormalisation}{n}$
\State $C\gets$ all boundary-aware phrase matches in $u$
\State Sort $C$ by start position and descending matched length
\State $M\gets$ empty accepted list; $O\gets$ empty occupied intervals
\For{candidate $c$ in $C$}
  \If{$c$ does not overlap an interval in $O$}
    \State Accept $c$ into $M$ and add its interval to $O$
  \EndIf
\EndFor
\State Sort $M$ by textual position
\For{mention $m$ in $M$}
  \State Inspect bounded context and positive-negation idiom exceptions
  \State Determine other-subject, historical, negated, uncertain or present status
  \State Record canonical phrase, matched text, status, reason and offsets
  \State Update the latest assertion for $m$'s canonical symptom
\EndFor
\State Group retained canonical symptoms by their latest assertion status
\State $b\gets$ comma-separated distinct present symptoms in first-mention order
\State \Return $(E,b)$
\end{algorithmic}
\end{algorithm}
No disease label is inferred here. Offsets refer to the normalised matching text.
The latest mention determines status, but arbitrary implicit corrections are not parsed.

\newpage
\begin{algorithm}[H]
\caption{Choose urgency by the highest-priority matching rule}
\label{alg:urgency}
\begin{algorithmic}[1]
\Require Narrative $n$, extracted evidence $E$, configured urgency rules $R$
\Ensure Urgency level $u$, supporting rules and classification-bypass decision
\State Normalise the narrative for rule matching
\State Collect asserted current symptoms from $E$
\State Collect eligible raw warning-phrase hits with assertion checks
\State Exclude clearly negated, historical and other-subject phrase hits
\State Allow specified uncertain raw warning phrases to remain actionable
\State $M\gets$ empty matched-rule list
\For{rule $r$ in $R$}
  \If{the implemented combination of evidence satisfies $r$}
    \State Append $r$'s identifier, urgency level and supporting evidence to $M$
  \EndIf
\EndFor
\State Rank levels as emergency $>$ urgent $>$ routine
\State $u\gets$ highest level in $M$, or routine when $M$ is empty
\State $bypass\gets(u\in\{\text{emergency},\text{urgent}\})$
\State Select the configured action message for $u$
\State \Return $(u,bypass,M,\text{action message})$
\end{algorithmic}
\end{algorithm}
Routine means no configured warning rule matched; it is not a clinical declaration
of safety. Emergency and urgent outcomes both take priority over disease prediction.

\newpage
\begin{algorithm}[H]
\caption{Run the trained BioBERT model and apply score display gates}
\label{alg:classify}
\begin{algorithmic}[1]
\Require Canonical current-symptom text $b$, local checkpoint $q$, thresholds
\Ensure Ranked model outputs $P$, status $s$, display flag $d$
\State Validate checkpoint files, 773 labels and classification-head tensors
\State Load and cache the trained model and tokeniser from $q$ locally
\State Tokenise $b$ with truncation at 128 tokens
\State Run the model in evaluation and inference-only mode
\State $p\gets\Call{Softmax}{\text{model logits}}$
\State $P\gets$ three highest scores and their checkpoint label names
\State Round recorded scores to six decimal places
\State $m\gets P_1.\text{score}-P_2.\text{score}$; record the rounded margin
\If{$P_1.\text{score}<0.50$}
  \State $s\gets\text{low score}$
\ElsIf{$m<0.10$}
  \State $s\gets\text{ambiguous prediction}$
\Else
  \State $s\gets\text{prediction available}$
\EndIf
\State $d\gets(s=\text{prediction available})$
\State \Return $(P,s,d)$
\end{algorithmic}
\end{algorithm}
The threshold values are configurable prototype heuristics, not calibrated disease
probabilities. The orchestration stage can still withhold an otherwise available result.

\newpage
\begin{algorithm}[H]
\caption{Orchestrate extraction, urgency, clarification and classification}
\label{alg:pipeline}
\begin{algorithmic}[1]
\Require Confirmed narrative $n$, extractor, rules and model runtime
\Ensure Analysis object with response, evidence, decisions and timings
\State Reject empty input and create a normalised text copy
\State $(E,b)\gets\Call{ExtractSymptoms}{n}$
\State $U\gets\Call{DetectUrgency}{n,E}$
\State $Q\gets\Call{GenerateClarification}{n,E,U,\text{no classification},1}$
\State $ran\gets\text{false}$
\If{$U.\text{bypass}$}
  \State $C\gets$ skipped classification with the urgency reason
\ElsIf{$Q.\text{pause classification}$}
  \State $C\gets$ skipped classification with the missing/uncertain symptom reason
\ElsIf{$b$ is non-empty}
  \State $ran\gets\text{true}$; $C\gets\Call{ClassifyWithGates}{b}$
  \State $Q\gets\Call{GenerateClarification}{n,E,U,C,1}$
  \If{$Q$ requires withholding}
    \State Hide the prediction and record low-score or ambiguity clarification status
  \EndIf
\Else
  \State $C\gets$ skipped classification with no-current-symptom reason
\EndIf
\State Store whether model inference actually occurred
\State $A\gets\Call{GenerateRuleResponse}{n,E,U,Q,C}$
\State Assemble request ID, inputs, evidence, stage trace and measured timings
\State \Return the analysis object including $A$
\end{algorithmic}
\end{algorithm}
The normalised text copy is retained for inspection. Extraction and urgency use
the full narrative, while the trained classifier receives the canonical symptom list.

\newpage
\begin{algorithm}[H]
\caption{Process one clarification answer without losing the original context}
\label{alg:followup}
\begin{algorithmic}[1]
\Require Original narrative $n$, free-form answer $a$
\Ensure Updated analysis or a closed-loop non-diagnostic fallback
\State $n'\gets\Call{Strip}{n}$
\State $a'\gets\Call{Strip}{a}$
\State $c\gets\Call{JoinWithSpace}{n',a'}$
\State $Z\gets\Call{AnalysePipeline}{c}$
\If{$Z.\text{response type}=\text{clarification}$}
  \State Preserve the current classification status as the pre-limit status
  \State Set classification status to clarification-limit-reached
  \State Hide the disease-class result
  \State Close clarification; remove questions; record the round-limit reason
  \State Set the clarification trace stage to limit-reached
  \State Regenerate the rule-based fallback with the updated decisions
\EndIf
\State Retain $n$, $a$ and $c$ as separate conversation fields
\State Mark the input source as clarification answer
\State \Return $Z$
\end{algorithmic}
\end{algorithm}
Emergency or urgent evidence in the answer is reconsidered immediately. A reply
such as no other symptoms does not erase a symptom recognised in the original narrative.

\newpage
\begin{algorithm}[H]
\caption{Select deterministic response content from pipeline state}
\label{alg:response}
\begin{algorithmic}[1]
\Require Current symptoms $E$, urgency $U$, clarification $Q$, classification $C$
\Ensure Structured response with actions, guidance, sources and disclaimer
\If{$U=\text{emergency}$}
  \State Select emergency action; hide class and omit routine self-care
\ElsIf{$U=\text{urgent}$}
  \State Select prompt-assessment action; hide class and omit routine self-care
\ElsIf{$Q$ requires clarification or withholding}
  \State Select clarification wording and the free-form question; hide the class
\ElsIf{$C.\text{display prediction}=\text{true}$}
  \State Describe the accepted output as a research dataset class, not a diagnosis
  \State Match symptom-level guidance using current symptoms in $E$
\Else
  \State Select a professional-assessment fallback without a disease label
  \If{the clarification round limit was reached}
    \State Explain insufficient symptoms or an inferred but withheld result
    \State Match any supported symptom-level guidance using $E$
  \EndIf
\EndIf
\State Deduplicate actions, guidance messages and source identifiers
\State Include applicable safety-net wording and the research disclaimer
\State Assemble structured fields and their text representation
\State \Return the response
\end{algorithmic}
\end{algorithm}
Guidance is symptom-based and does not implement treatment for 773 disease
classes. Source-linked content still requires clinical review of its applicability.

\newpage
\begin{algorithm}[H]
\caption{Evaluate the legacy paired speech and noise experiment}
\label{alg:legacy-eval}
\begin{algorithmic}[1]
\Require Selected recordings $S$, references, one noise file, fixed models, seed
\Ensure Per-condition records and descriptive summaries
\State Load the fixed Whisper recogniser and trained BioBERT classifier
\For{sample $s$ in $S$}
  \State Obtain a direct BioBERT prediction from the reference transcript
  \For{condition $c\in\{\text{clean},20,10,5,0\text{ dB}\}$}
    \State Create the condition waveform using a reproducible noise segment
    \For{route $r\in\{\text{original},\text{enhanced}\}$}
      \State Apply enhancement only when $r=\text{enhanced}$
      \State $h\gets\Call{WhisperTranscribe}{\text{route waveform}}$
      \State Normalise reference and $h$ with the same error-rate procedure
      \If{the normalised hypothesis is non-empty}
        \State Compute WER and CER; classify $h$ directly with BioBERT
      \Else
        \State Record WER and CER as 1.0 and an empty disease prediction
      \EndIf
      \State Record prediction agreement with the reference-text prediction
      \State Record disease correctness only if an appropriate mapped label exists
      \State Record score, low-score flag and ranked classes
      \State Save the accumulated per-condition records incrementally
    \EndFor
  \EndFor
\EndFor
\State Group by dataset, condition and route; calculate descriptive means
\State \Return records and summaries
\end{algorithmic}
\end{algorithm}
This implemented pilot contains six selected recordings and sixty repeated
condition/route observations. It bypasses the current extractor and conversation
controller, so it is not the completed end-to-end evaluation of the latest prototype.
\end{document}
